\subsection{TOTAL ALGEBRA OF A MONOID}

The algebra of a monoid S over A is (as an A-module) the submodule of the product \(A^{s}\) consisting of the families \((\alpha_{s})_{s\in\mathbb{S}}\) of finite support; the multiplication in this algebra is defined by the relations \((\alpha_{s})(\beta_{s})=(\gamma_{s})\) , where, for all \(s\in S\) ,

\[
\gamma_ {s} = \sum_ {t u = s} \alpha_ {t} \beta_ {u}\tag{38}
\]

(cf.~no. 6, formula (35)). The sum on the right hand side of (38) is meaningful because \((\alpha_{s})\) and \((\beta_{s})\) are families of finite support and so therefore is the double family \((\alpha_{t}\beta_{u})_{(t,u)\in S\times S}\) . But the right hand side of (38) is also meaningful for arbitrary elements \((\alpha_{s})\) , \(\beta_{s})\) of \(A^{S}\) when the monoid \(S\) satisfies the following condition:

\begin{enumerate}
\def\labelenumi{(\Alph{enumi})}
\setcounter{enumi}{3}
\tightlist
\item
  For all \(s \in \mathbf{S}\) , there exists only a finite number of ordered pairs \((t, u)\) in \(\mathbf{S} \times \mathbf{S}\) such that \(tu = s\) .
\end{enumerate}

Suppose then that S satisfies condition (D); a multiplication law can then be defined on the product A-module \(A^{s}\) by formula (38). It is immediate that the multiplication thus defined on \(A^{s}\) is A-bilinear; also it is associative, since, for \(\alpha\) , \(\beta\) , \(\gamma\) in \(A^{s}\) ,

\[
\sum_ {u v w = t} \alpha_ {u} \beta_ {v} \gamma_ {w} = \sum_ {r w = t} \left(\left(\sum_ {u v = r} \alpha_ {u} \beta_ {v}\right) \gamma_ {w}\right) = \sum_ {u s = t} \left(\alpha_ {u} \left(\sum_ {v w = s} \beta_ {v} \gamma_ {w}\right)\right).
\]

This multiplication and the A-module structure on \(A^{s}\) therefore define on \(A^{s}\) a unital associative algebra structure over A; we shall say that the set \(A^{s}\) , with this structure, is the total algebra of the monoid S over A.

It is immediate that the algebra \(\mathbf{A}^{(\mathbf{S})}\) of the monoid \(\mathbf{S}\) over \(\mathbf{A}\) (also called the restricted algebra of \(\mathbf{S}\) when necessary to avoid confusion) is a subalgebra of the total algebra of \(\mathbf{S}\) over \(\mathbf{A}\) (and is identical with the latter when \(\mathbf{S}\) is finite). By an abuse of language, every element \((\xi_s)_{s \in \mathbf{S}}\) of the total algebra of \(\mathbf{S}\) over \(\mathbf{A}\) is also denoted by the same notation \(\sum_{s \in \mathbf{S}} \xi_s e_s\) (or even \(\sum_{s \in \mathbf{S}} \xi_s s\) ) as the elements of the restricted algebra of \(\mathbf{S}\) ; of course the summation symbol appearing in this notation corresponds to no algebraic operation because it is taken over an infinity of terms \(\neq 0\) in general. With this notation, multiplication in the total algebra of \(\mathbf{S}\) is also given by formula (35) of no. 6.

If S is commutative, so is its total algebra \(A^{S}\) . If T is a submonoid of S, the total algebra \(A^{T}\) of the monoid is canonically identified with a subalgebra of the total algebra of S. If \(\rho: A \to B\) is a ring homomorphism, the canonical extension \(\rho^{S}: A^{S} \to B^{S}\) is an A-homomorphism of the total algebra of S over A into the total algebra of S over B, which extends the canonical homomorphism \(A^{(S)} \to B^{(S)}\) .
% semantic-fragments: legacy-input-seam
\relax{}
